\documentclass[11pt]{article}
\usepackage[margin=1in]{geometry}
\usepackage{amsmath,amssymb}
\usepackage{enumitem}
\setlength{\parindent}{0pt}
\setlength{\parskip}{6pt}
\begin{document}
\section*{STAT 412 HW07 --- Question 6}

\subsection*{Problem}

\begin{quote}
The average height of females in the freshman class of a certain college has historically been 162.8 centimeters with a standard deviation of 6.8 centimeters. Is there reason to believe that there has been a change in the average height if a random sample of 48 females in the present freshman class has an average height of 165.8 centimeters? Use a P-value in your conclusion. Assume the standard deviation remains the same.
\end{quote}

\subsection*{Solution}

The word ``changed'' makes this a two-tailed test:
\[
H_0:\mu=162.8,\qquad H_1:\mu\ne162.8.
\]

Since the population standard deviation is assumed known, use a $z$ statistic:
\[
z=\frac{\bar{x}-\mu_0}{\sigma/\sqrt{n}}
  =\frac{165.8-162.8}{6.8/\sqrt{48}}
  \approx 3.0566.
\]

Rounded to two decimals, the test statistic is
\[
z=3.06.
\]

The P-value is two-tailed:
\[
P\text{-value}=2P(Z>3.0566)\approx 0.0022389.
\]

The platform prompt says to round the P-value to three decimal places, so enter
\[
0.002.
\]
The value $0.0022$ is the four-decimal version, but it is not rounded to the requested three decimals.

\subsection*{Final Answer}

\fbox{\begin{minipage}{0.94\linewidth}
$H_0:\mu=162.8$; $H_1:\mu\ne162.8$. Test statistic $z=3.06$. P-value $=0.002$ to three decimals. Reject $H_0$; there is sufficient evidence that the mean height has changed.
\end{minipage}}

\end{document}
